% ══════════════════════════════════════════════════════════════
%  Chapter 4: Bitcoin Whitepaper in Rust
% ══════════════════════════════════════════════════════════════
\chapter{Bitcoin Whitepaper in Rust}
\label{ch:whitepaper-rust}

\begin{learnbox}
\begin{itemize}
  \item Map each Bitcoin whitepaper concept to a concrete Rust data structure
  \item Explain the hashing and serialization rules that connect whitepaper theory to working code
  \item Trace how each whitepaper section corresponds to modules in the repository
  \item Understand the business objects that form the foundation of the implementation
\end{itemize}
\end{learnbox}

\lettrine{N}{ow} that we've walked through the Bitcoin \index{whitepaper}whitepaper, we switch from \textbf{theory to implementation}. In the sections that follow, we translate the \index{whitepaper}whitepaper's ideas into \textbf{concrete \index{Rust!modules}Rust modules and methods}, with complete code listings and explanations---calling out what we implement exactly, what we intentionally simplify, and where real-world Bitcoin Core adds additional rules.

\begin{notebox}
\textbf{Tip:} Read this chapter alongside Chapter~4 (Bitcoin Whitepaper Summary). Each section here maps directly to a whitepaper section, and understanding the theory first makes the Rust encoding decisions much clearer.
\end{notebox}

% ──────────────────────────────────────────────────────────────
\section{Rust Refresher (Optional)}
\label{sec:rust-refresher}
\index{Rust!introduction}

If we want to brush up on Rust before working through the implementations, we read the Rust guide:

\begin{itemize}
  \item \textbf{Chapter~\ref{ch:rust-guide}: Rust Language Guide} --- start with \textbf{Rust Installation \& Setup}, then continue to \textbf{Introduction}.
\end{itemize}

% ──────────────────────────────────────────────────────────────
\section{Read Options}
\label{sec:read-options}

\textbf{Split sections (recommended)}: use the links below.

% ──────────────────────────────────────────────────────────────
\section{Split Sections (0--12 Plus Appendices)}
\label{sec:split-sections}

The material in this chapter is organized as follows:

\begin{itemize}
  \item Business objects
  \item Introduction --- Bitcoin Whitepaper Section 1
  \item Transactions chain of signatures --- Bitcoin Whitepaper Section 2
  \item Timestamp server block header chaining --- Bitcoin Whitepaper Section 3
  \item Proof of work --- Bitcoin Whitepaper Section 4
  \item Network operation --- Bitcoin Whitepaper Section 5
  \item Incentive mechanism --- Bitcoin Whitepaper Section 6
  \item Reclaiming disk space --- Bitcoin Whitepaper Section 7
  \item Merkle trees and SPV --- Bitcoin Whitepaper Sections 7--8
  \item Combining splitting value --- Bitcoin Whitepaper Section 9
  \item Privacy --- Bitcoin Whitepaper Section 10
  \item Confirmations and attacker probability --- Bitcoin Whitepaper Section 11
  \item Conclusion --- Bitcoin Whitepaper Section 12
  \item Appendix A: Object connectivity end to end flow
  \item Appendix B: Mapping to this repository
\end{itemize}

% ──────────────────────────────────────────────────────────────
\section{How to Read This Section (Node Lifecycle)}
\label{sec:node-lifecycle}

We structure this section around the core node pipeline, from parsing bytes to committing state. In Rust terms, the data path is:

\[\text{\textbf{bytes} \(\to\) \textbf{structs} \(\to\) \textbf{canonical bytes} \(\to\) \textbf{hashes} \(\to\) \textbf{validation} \(\to\) \textbf{state update}}\]

\begin{itemize}
  \item \textbf{We accept a \index{transaction}transaction}: parse bytes into \code{\index{transaction}Transaction} \(\to\) serialize canonically \(\to\) compute \code{\index{txid}txid} \(\to\) validate (authorization + value rules) \(\to\) check inputs are unspent in our \index{UTXO}UTXO view \(\to\) add to \index{mempool}mempool \(\to\) relay.
  \item \textbf{We accept a \index{block}block}: parse into \code{\index{block}Block} \(\to\) validate \code{\index{block header}BlockHeader} (prev link + PoW target) \(\to\) validate tx set (including ``not already spent'') \(\to\) apply \index{UTXO}UTXO transition atomically \(\to\) maybe reorg \(\to\) relay.
\end{itemize}

When we feel lost, we come back to this Rust checklist (whitepaper phrase \(\to\) concrete implementation question):

\begin{itemize}
  \item \textbf{Whitepaper says ``hash''}: ``What exact bytes are hashed, in what order, and with what endianness?'' (consensus serialization)
  \item \textbf{Whitepaper says ``sign''}: ``What exact digest do we sign/verify per-input, and what prevout/script context is included?'' (signature hashing / script context)
  \item \textbf{Whitepaper says ``not already spent''}: ``Does \code{TxIn.previous\_output: OutPoint} exist in our UTXO view at this tip?'' (state + double-spend prevention)
\end{itemize}

% ──────────────────────────────────────────────────────────────
\section{Implementation Roadmap (With Rust Encoding)}
\label{sec:implementation-roadmap}

\begin{steps}
\step{Bytes (canonical encoding)} Implement \index{CompactSize}CompactSize/varints and deterministic \index{serialization}serialization for the core objects. Outcome: ``same object \(\to\) same bytes'' (always).

\step{Identity (hashes as IDs)} Implement \code{\index{SHA-256}sha256d}, \code{\index{txid}txid}, \index{Merkle tree}Merkle root, and \code{\index{block hash}block\_hash}. Outcome: ``same bytes \(\to\) same IDs'' across nodes.

\step{Authorization (spends are permitted)} Implement \index{script}script/\index{digital signature}signature verification using the referenced output as the spend context. Outcome: ``only the key-holder can spend'' (no forged spends).

\step{State (UTXO transitions)} Implement \index{UTXO}UTXO updates as a pure state transition: spend inputs, create outputs. Outcome: ``not already spent'' is enforceable at the state layer.

\step{Consensus (block acceptance + best chain)} Implement \index{block!validation}block validation and \index{chain selection}chain selection (and \index{reorg}reorg handling if needed). Outcome: nodes converge on the same history under the \index{longest chain rule}longest/most-work rule.
\end{steps}

% ──────────────────────────────────────────────────────────────
\section{What the Whitepaper Specifies (and What It Doesn't)}
\label{sec:whitepaper-scope}

In the summary chapter we focused on \emph{what} the paper says. Here we focus on \emph{how that becomes code} without losing the paper's intent and meaning.

The whitepaper specifies the \emph{core relationships}:

\begin{itemize}
  \item \textbf{Transactions} are transfers of ownership via digital signatures (Bitcoin Whitepaper Section 2).
  \item \textbf{Blocks} timestamp transactions by hashing them into a chain (Bitcoin Whitepaper Section 3).
  \item \textbf{Proof-of-work} makes rewriting history computationally impractical (Bitcoin Whitepaper Section 4).
  \item \textbf{Nodes} broadcast transactions and blocks and converge on the ``longest'' chain (Bitcoin Whitepaper Section 5).
  \item \textbf{Merkle trees} enable compact commitments and SPV proofs (Bitcoin Whitepaper Sections 7--8).
\end{itemize}

The deployed Bitcoin protocol adds important concrete details beyond the paper (still consistent with it), such as:

\begin{itemize}
  \item \textbf{double-SHA256} (``hash twice'') for txids and block hashes
  \item \textbf{CompactSize/varints} and strict little-endian \textbf{consensus serialization}
  \item \textbf{Script} (e.g., P2PKH/P2WPKH; see description in next section) and signature hashing (``sighash'') rules
  \item a real P2P message protocol (inv/getdata, headers, etc.)
\end{itemize}

The whitepaper itself does \textbf{not} fully specify:

\begin{itemize}
  \item exact wire encodings (byte-for-byte consensus serialization)
  \item Script details, signature hashing rules, and full validation edge cases
  \item production P2P protocol message formats
\end{itemize}

That's fine for learning: the whitepaper is the conceptual contract, and implementations supply the missing exactness. Our goal is to make the missing pieces explicit (especially byte encoding), because consensus is ultimately ``agree on bytes''.

% ──────────────────────────────────────────────────────────────
\section{Script: P2PKH vs P2WPKH}
\label{sec:p2pkh-p2wpkh}
\index{script}\index{P2PKH}\index{P2WPKH}

In Bitcoin, most outputs are ``locked'' by a \textbf{\index{scriptPubKey}scriptPubKey} (a small stack program). Spending an output means providing data that makes that program succeed. Two common standard output types are:

\begin{deepdive}{P2PKH (Pay-to-PubKey-Hash)}

\textbf{\index{P2PKH}P2PKH (Pay-to-PubKey-Hash)}: the output commits to a 20-byte \code{\index{HASH160}HASH160(pubkey)} and can be spent by providing a \index{digital signature}signature and the corresponding \index{public key}public key.

\textbf{Locking form (scriptPubKey)}:
\begin{minted}[fontsize=\footnotesize,linenos=false]{text}
OP_DUP OP_HASH160 <20-byte pubKeyHash> OP_EQUALVERIFY OP_CHECKSIG
\end{minted}

\textbf{Spending data (conceptually)}: \code{<sig> <pubkey>} (the node hashes \code{pubkey} to check it matches \code{pubKeyHash}, then verifies \code{sig} under the signature-hash rules).

\textbf{Common address encoding}: legacy Base58 addresses that start with \code{1...} on mainnet.

\end{deepdive}

\begin{infobox}{Why \code{HASH160(pubkey)} and not just \code{SHA256(pubkey)}?}

In standard \textbf{\index{script}Script/address constructions} (like \index{P2PKH}P2PKH/\index{P2WPKH}P2WPKH), Bitcoin uses \code{\index{HASH160}HASH160 = \index{RIPEMD160}RIPEMD160(\index{SHA-256}SHA256(x))} for pubkey hashes largely because:

\begin{itemize}
  \item It is \textbf{smaller} (20 bytes vs 32 bytes), which reduces output sizes and UTXO set footprint.
  \item It provides \textbf{defense in depth} by combining two hash functions (a practical break would need to affect both).
  \item It is a \textbf{historical standard} baked into early script/address conventions (changing it would be a broad compatibility and consensus break).
\end{itemize}

Note: this is separate from block/tx identifiers. \textbf{txids} (and \textbf{block hashes}) are computed with \textbf{double-SHA256} (i.e., \code{SHA256(SHA256(bytes))}) over consensus-serialized data (tx serialization for txids; block header serialization for block hashes).

\end{infobox}

\begin{deepdive}{P2WPKH (Pay-to-Witness-PubKey-Hash, SegWit v0)}

\textbf{\index{P2WPKH}P2WPKH (Pay-to-Witness-PubKey-Hash, \index{SegWit}SegWit v0)}: the same ownership model as \index{P2PKH}P2PKH (pubkey hash), but \index{digital signature}signatures and pubkeys are carried in \textbf{\index{witness}witness} data (\index{SegWit}SegWit), not the legacy \index{scriptSig}scriptSig path.

\textbf{Locking form (witness program)}:
\begin{minted}[fontsize=\footnotesize,linenos=false]{text}
0 <20-byte pubKeyHash>
\end{minted}

\textbf{Spending data}: witness stack provides \code{<sig> <pubkey>} (native P2WPKH has an empty \code{scriptSig}).

\textbf{Why it matters}: witness data is excluded from the legacy txid, which eliminates third-party malleability for signature data and makes blocks more space-efficient under weight rules.

\textbf{Common address encoding}: Bech32 addresses that start with \code{bc1q...} on mainnet. (For compatibility there is also P2SH-wrapped P2WPKH, often starting with \code{3...}.)

\end{deepdive}

% ──────────────────────────────────────────────────────────────
\section{Further Reading}
\label{sec:ch04-further-reading}

\begin{itemize}
  \item \textbf{Bitcoin Whitepaper} --- Read alongside this chapter for the full context behind each Rust encoding. See \url{https://bitcoin.org/bitcoin.pdf}
  \item \textbf{serde Documentation} --- The serialization framework used throughout the implementation. See \url{https://serde.rs/}
  \item \textbf{Bitcoin Developer Reference} --- Technical specifications for Bitcoin data structures. See \url{https://developer.bitcoin.org/reference/}
\end{itemize}

% ──────────────────────────────────────────────────────────────
\begin{chaptersummary}
  \item We mapped every Bitcoin whitepaper concept to a concrete Rust data structure, showing how theory translates into working code.
  \item We defined the core business objects --- Block, Transaction, BlockHeader, and their supporting types --- that form the implementation's foundation.
  \item We traced hashing and serialization rules from whitepaper descriptions to Rust implementations using SHA-256 and serde.
  \item We established the mapping between whitepaper sections and repository modules, creating a bridge between concept and code.
  \item In the next chapter, we examine how the Rust project is organized --- the crate workspace, dependency graph, and build configuration that make all of this compile and run.
\end{chaptersummary}
